% ==================================================================
\section{Chapter 3 --- Filtrations and Martingales}
% ==================================================================
Throughout this chapter we consider a probability space $(\Omega, \mc{F}, P)$.
\subsection{Filtrations and processes}
\begin{definition}[Filtration and filtered probability space]\label{chap3.1:definition-filtration}
\LG{Definition 3.1} \\
A \textit{filtration} on $(\Omega, \mc{F}, P)$ is a collection $(\mc{F}_t)_{0 \leq t \leq \infty}$ of sub-$\sigma$-algebras of $\mc{F}$ indexed by $[0,\infty]$, such that $\mc{F}_s \subseteq \mc{F}_t$ for every $0 \leq s \leq t \leq \infty$. We thus have, for every $0 \leq s < t$ $$\mc{F}_0 \subseteq \mc{F}_s \subseteq \mc{F}_t \subseteq \mc{F}_\infty \subseteq \mc{F}.$$
We also say that $(\Omega, \mc{F}, (\mc{F}_t), P)$ is a \textit{filtered probability space}.
\end{definition}

\begin{definition}[Canonical filtration]\label{chap3.1:definition-canonical_filtration}
\LG{p. 42} \\
Let $X = (X_t)_{t \geq 0}$ be any random process indexed by $\R_+$. The \textit{canonical filtration of $X$}  is defined by
$$\mc{F}_t^X:=\sigma(X_s, 0 \leq s \leq t)=\sigma\left(\bigcup_{s \leq t}\sigma(X_s)\right), \quad \mc{F}_\infty^X=\sigma(X_s, s \geq 0)=\sigma\left(\bigcup_{s \geq 0}\sigma(X_s)\right)$$
\end{definition}

\begin{definition}[Adaptedness]\label{chap3.1:definition-adaptedness}
\LG{Definition 3.3} \\
A random process $X=(X_t, t \geq 0)$ is called \textit{adapted} to a filtration $(\mc{F}_t)_{t \geq 0}$ if, for every $t \geq 0$, then $X_t$ is $\mc{F}_t$-measurable.
\end{definition}

Clearly $X$ is adapted to its canonical filtration by construction, which the following proposition makes clear.
\begin{proposition}\label{chap3.1:proposition-canonical_filtration_is_smallest_filtration_for_adaptness}
$(F^X_t)_{t \geq 0}$ is the smallest filtration to which $X$ is adapted. That is if $\mc{G}$ is any filtration that $X$ is adapted to, then $\mc{F}_t^X \subseteq \mc{G}$.
\end{proposition}
\begin{proof}
\textbf{Filtration property:} Let $s \leq t$. Then note
$$\bigcup_{u\leq s}\sigma(X_u) \subseteq\bigcup_{u \leq t} \sigma(X_u)$$
so every generator $X_u$ of $\mc{F}_s^X$ satisfies $u \leq s \leq t$, so it is also a generator of $\mc{F}_s^X$, which gives $\mc{F}_s^X \subseteq \mc{F}_t^X$. Also since $X_u$ is $\mc{F}$-measurable then $\mc{F}_t^X \subseteq \mc{F}$.\\
\textbf{Adaptedness:}. For each $t$, then $\sigma(X_t) \subseteq\mc{F}_t^X$ by construction, so therefore $X_t$ is $\mc{F}_t^X$-measurable.\\
\textbf{Minimality:} Let $(\mc{G}_t)$ be a filtration with $X$ adapted to it. Fix $t \geq 0$. For any $s \leq t$:
\begin{itemize}
    \item $X_s$ is $\mc{G}_s$-measurable, because $X$ is adapted to $(\mc{G}_t)$.
    \item $\mc{G}_s \subseteq \mc{G}_t$, because $(\mc{G}_t)$ is a filtration.
\end{itemize}
So $\sigma(X_s) \subseteq \mc{G}_t$ for every $s \leq t$. This means $\bigcup_{s \leq t}\sigma(X_s) \subseteq \mc{G}_t$. To see this last point then let $A \in \bigcup_{s \leq t}\sigma(X_s)$. Then there exists $0 \leq s' \leq t$ such that $A \in \sigma(X_{s'}) \subseteq\mc{G}$, which shows the containment.\\
Since $\mc{G}_t$ is a $\sigma$
-algebra, it also contains the smallest $\sigma$-algebra generated by that union that is
$$\mc{F}_t^X \subseteq \mc{G}_t$$
This shows the wanted minimality.
\end{proof}

\begin{definition}[Right-continuous filtration]\label{chap3.1:definition-right_cont_filtration} 
\LG{p. 42} \\
Let $(\mc{F})_{0 \leq t \leq \infty}$ be a filtration on $(\Omega, \mc{F}, P)$. For every $t \geq 0$, set
$$\mc{F}_{t+}=\bigcap_{s>t}\mc{F}_s$$
and $\mc{F}_{\infty+}=\mc{F}_\infty$. Note that $\mc{F}_t \subseteq \mc{F}_{t+}$ and that the collection $(\mc{F}_{t+})_{0 \leq t \leq \infty}$ is also a filtration. We say that the filtration $(\mc{F}_t)$ is \textit{right-continous} if
$$\mc{F}_{t+}=\mc{F}_t, \quad \forall t \geq 0.$$
By construction, the filtration $(\mc{F}_{t+})$ is right-continous.
\end{definition}

\begin{definition}[Completeness of filtrations]\label{chap3.1:definition-completeness_filtration}
\LG{p. 42} \\
Let $(\mc{F}_t)$ be a filtration and let $\mc{N}$ be the class of all $(\mc{F}_\infty,P)$-negligible sets, that is
$$A \in \mc{N} \quad \Leftrightarrow \quad \exists A' \in \mc{F}_\infty \text{ such that } A \subseteq A' \text{ and } P(A')=0$$
The filtration is said to be \textit{complete} if $\mc{N} \subseteq \mc{F}_0$ and thus $\mc{N} \subseteq \mc{F}_t$ for every $t \in [0,\infty]$.
\end{definition}
We will usually assume that the filtration we are working with is \textit{complete and right-continous}. This will ensure path regularity, and make stopping times and martingales well-behaved.
Without the completeness-assumptions we would not be able to infer that quantities like $\sup_{s \leq t}X_s$ are measurable, and one might end up with an event that is not contained in the $\sigma$-algebra.

Likewise right-continuity of filtrations ensures that hitting times of open sets is a stopping time, because such a hitting time requires looking an infinitesimal amount of time ahead. 

\textbf{One can always consider the completion of the filtration, and thereby safely avoid issues related to nasty null-sets in the following way:}
\begin{remark}[Completing a non-complete filtration] \label{chap3.1:remark-completing_a_non_complete_filtration}
\LG{p. 42} \\
For a given filtraation $(\mc{F}_t)_{t \geq 0}$ it possible to consider its \textit{completion} $\mc{F}'_t=\sigma\left(\mc{N} \cup \mc{F}_t\right)$, where $\mc{N}$ again is the collection of negligible elements of $\mc{F}_\infty$. Often we would apply this completion procedure to the canonical filtration of a random process, and call the resulting filtration the \textit{completed canonical filtration}.
\end{remark}

\begin{trap}\label{chap3.1:trap-right_continous_process_does_not_imply_right_continous_filtration}
The canonical filtration of a right-continuous process is \textbf{not right-continuous} in general. A counter example found in \cite{chafai2021completeness} is given by $X_t=tZ$ for all $t \geq 0$ where $Z$ is a non-constant random variable. Indeed, we have $\sigma(X_0)=\{\emptyset, \Omega\}$ while $\sigma(X_{0+\epsilon}: \epsilon > 0) = \sigma(Z) \neq \sigma(X_0)$. However it can be shown that the completion of the canonical filtration of a Feller-Markov process, including all Lévy processes and in particular Brownian motion, is always right-continuous.
\end{trap}


\begin{definition}[Measurable and progressive processes] \label{chap3.1:definition-measurable_and_progressive_processes}
\LG{Definitions 3.2 and 3.3}
A random process $(X_t)_{t \geq 0}$ with values in a measurable space
$(E,\mc{E})$ is said to be \textit{measurable} if the mapping

$$(\omega,s) \mapsto X_s(\omega)$$

defined on $\Omega \times \R_+$ equipped with the product
$\sigma$-algebra $\mc{F} \otimes \mc{B}(\R_+)$ is measurable.

A process is said to be \textit{progressive} or
\textit{progressively measurable} if, for every $t \geq 0$, the mapping

$$(\omega,s) \mapsto X_s(\omega)$$

defined on $\Omega \times [0,t]$ is measurable with respect to
$\mc{F}_t \otimes \mc{B}([0,t])$.
\end{definition}
\begin{proposition}[Local characterization of measurability]
\LG{p. 43, stated in first paragraph.} \\
A process $(X_t)_{t\geq0}$ is measurable if and only if, for every
$t\geq0$, the mapping
\[
(\omega,s)\mapsto X_s(\omega)
\]
restricted to $\Omega\times[0,t]$ is measurable with respect to
\[
\mc{F}\otimes\mc{B}([0,t]).
\]
\end{proposition}

\begin{proof}
The forward implication is immediate. Suppose that
$$(\omega,s)\mapsto X_s(\omega)$$
is measurable on $\Omega\times\mathbb R_+$ with respect to
$\mathcal F\otimes\mathcal B(\mathbb R_+)$. For any $t\geq0$, the set
$\Omega\times[0,t]$ is measurable, since $[0,t]\in\mathcal B(\mathbb R_+)$.
Hence, restricting the mapping to $\Omega\times[0,t]$ gives a measurable
mapping with respect to
$$\mathcal F\otimes\mathcal B([0,t]).$$
Conversely, suppose that for every $t\geq0$ the restricted mapping
$$(\omega,s)\mapsto X_s(\omega),
\qquad (\omega,s)\in\Omega\times[0,t],$$
is $\mc{F}\otimes\mc{B}([0,t])$-measurable.
Let $A\in\mc{E}$. For every $n\in\N$, we then have
$$\{(\omega,s)\in\Omega\times[0,n]:X_s(\omega)\in A\}
\in
\mc{F}\otimes\mc{B}([0,n]).$$
Since
$$\mc{F}\otimes\mc{B}([0,n])
\subseteq
\mc{F}\otimes\mc{B}(\R_+),$$
each of these sets belongs to $\mc{F}\otimes\mc{B}(\R_+)$.
Moreover,
$$\{(\omega,s)\in\Omega\times\R_+:X_s(\omega)\in A\}
=
\bigcup_{n=1}^{\infty}
\{(\omega,s)\in\Omega\times[0,n]:X_s(\omega)\in A\}.$$
Since $\mc{F}\otimes\mc{B}(\R_+)$ is a $\sigma$-algebra, the right-hand
side belongs to $\mc{F}\otimes\mc{B}(\R_+)$.

Thus $(\omega,s)\mapsto X_s(\omega)$ is measurable.
\end{proof}

\begin{remark}[Progressive $\subseteq$ adapted + measurable] \label{chap3.1:remark-progressive_implies_adapted_plus_measurable}
A progressive process is both adapted and measurable.
Indeed, fix $t\geq0$. By progressive measurability, the mapping
$$
F_t:\Omega\times[0,t]\to E,
\qquad
F_t(\omega,s)=X_s(\omega),
$$
is $\mathcal F_t\otimes\mathcal B([0,t])$-measurable. Consider the mapping
$$
i_t:\Omega\to\Omega\times[0,t],
\qquad
i_t(\omega)=(\omega,t).
$$
This mapping is $\mathcal F_t$-measurable, and
$$
X_t=F_t\circ i_t.
$$
Hence $X_t$ is $\mathcal F_t$-measurable. Since $t$ was arbitrary, $(X_t)_{t\geq0}$ is adapted. Moreover, since $\mathcal F_t\subseteq\mathcal F$,
$$
\mathcal F_t\otimes\mathcal B([0,t])
\subseteq
\mathcal F\otimes\mathcal B([0,t]).
$$
Thus the process satisfies the local characterization of measurability, and is therefore measurable.
\end{remark}


\begin{proposition} \label{chap3.1:proposition-right_continuity_implies_progressivity}
\LG{Proposition 3.4}
Let $(X_t)_{t \geq 0}$ be a random process with values in a metric space $(E, \mc{E})$ (equipped with its Borel $\sigma$-algebra). Suppose that $X$ is adapted and that the sample paths of $X$ are right-continous. That is for every $\omega \in \Omega, t \mapsto X_t(\omega)$ is right-continous. Then $X$ is progressive and this holds if one replaces right-continuity by left continuity.
\end{proposition}
\begin{proofidea}
The right-continuous version:
\begin{enumerate}
    \item For every $n \geq 1$ and $s \in [0,t]$ define $X_s^n=\begin{cases}X_{kt/n}, \quad &s \in [(k-1)t/n, kt/n), k \in \{1,...,n\} \\ X_t, \quad &s = t \end{cases}$
    \item Note that right continuity of sample paths ensure that $X_s(\omega)=\lim_{n \to \infty} X_s^n(\omega)$ for every $s \in [0,t], \omega \in \Omega$.
    \item Show that $X_s^n(\omega)$ is $\mc{F}_t \otimes \mc{B}([0,t])$-measurable.
    \item Use that a pointwise limit of measurable functions is also measurable with the same measurability condition to conclude that $X$ is progressive.
\end{enumerate}
\end{proofidea}
\begin{remark}\label{chap3.1:remark-right_continuity_implies_progressivity_implies_adapted_plus_measurable}
Note that Proposition \ref{chap3.1:proposition-right_continuity_implies_progressivity} and Remark \ref{chap3.1:remark-progressive_implies_adapted_plus_measurable} shows that (under weak assumption on the metric space that $X$ attains values in)
$$\text{Adapted + right-continous sample paths} \implies \text{Progressive} \implies \text{Measurable and adapted}$$
\end{remark}

\begin{quizbox}
The implication chain $\text{Adapted}+\text{right- (or left-) continuous paths} \Rightarrow \text{Progressive} \Rightarrow \text{Adapted}+\text{Measurable}$ is the single most-reused fact of this subsection: every process we actually work with (BM, càdlàg martingales, stochastic-integral integrands) is progressive \emph{because} it is adapted with one-sided-continuous paths, never by checking the definition directly. \\
Also know \emph{why} we assume the \textit{usual conditions} (complete and right-continuous filtration): completeness makes running suprema like $\sup_{s\le t}X_s$ genuinely $\mc F_t$-measurable (needed for Doob's inequalities and OST) even though the sup is over an uncountable set, while right-continuity of $(\mc F_t)$ is exactly what makes hitting times of \emph{open} sets stopping times without needing the $\mc F_{t+}$ workaround (\cref{chap3.2:proposition-hitting_times} below).
\end{quizbox}

\begin{proposition}[The progressive $\sigma$-algebra]
\LG{p. 43} \\
The collection $\mc{P}$ of all sets $A \in \mc{F} \otimes \mc{B}(\R_+)$ such that the process $X_t(\omega)=1_A(\omega,t)$ is progressive forms a $\sigma$-algebra on $\Omega \times \R_+$, which is called \textit{the progressive $\sigma$-algebra}. \\
A subset $A \subseteq \Omega \times \R_+$ belongs to $\mc{P}$ if and only if, for every $t \geq 0$, then $A \cap (\Omega \times [0,t])$ belongs to $\mc{F}_t \otimes \mc{B}([0,t])$.
\end{proposition}

\begin{proposition} \label{chap3.1:proposition-progressivity_from_sigma_algebra_characterization}
\LG{p. 43, stated in the bottom} \\
A process $(X_t)_{t \geq 0}$ is progressive if and only if the mapping 
$$(\omega,t) \mapsto X_t(\omega)$$
is measurable on $\Omega \times \R_0$ equipped with the $\sigma$-algebra $\mc{P}$.
\end{proposition}
\begin{proof}
Suppose first that $(X_t)_{t\geq0}$ is progressive. Let
$B\in\mc E$. For every $t\geq0$, progressive measurability implies that
the mapping
$$
(\omega,s)\mapsto X_s(\omega)
$$
restricted to $\Omega\times[0,t]$ is
$\mc F_t\otimes\mc B([0,t])$-measurable. Hence
$$
\begin{aligned}
X^{-1}(B)\cap(\Omega\times[0,t])
&=
\{(\omega,s)\in\Omega\times[0,t]:X_s(\omega)\in B\}\\
&\in
\mc F_t\otimes\mc B([0,t]).
\end{aligned}
$$
Since this holds for every $t\geq0$, the characterization of the
progressive $\sigma$-algebra gives
$$
X^{-1}(B)\in\mc P.
$$
Thus $X$ is $\mc P$-measurable.

Conversely, suppose that $X$ is $\mc P$-measurable. Let $t\geq0$ and
$B\in\mc E$. Then
$$
X^{-1}(B)\in\mc P.
$$
By the characterization of $\mc P$,
$$
X^{-1}(B)\cap(\Omega\times[0,t])
\in
\mc F_t\otimes\mc B([0,t]).
$$
But
$$
X^{-1}(B)\cap(\Omega\times[0,t])
=
(X|_{\Omega\times[0,t]})^{-1}(B).
$$
Therefore the restriction
$$
X|_{\Omega\times[0,t]}
:
\Omega\times[0,t]\to E
$$
is $\mc F_t\otimes\mc B([0,t])$-measurable. Since $t$ was arbitrary,
$(X_t)_{t\geq0}$ is progressive.
\end{proof}

\subsubsection{Measurability of running supremum from completeness}
% Requires: amsmath, amssymb, amsthm with proposition/remark environments
This subsection is adapted from an almost identical theorem from \cite{chafai2021completeness}.
\begin{proposition}[Measurability of the running supremum]
Let $(\Omega,\mathcal F,\mathbb P)$ be a probability space and let
$(X_t)_{t\in\mathbb R_+}$ be a process with values in a topological space $E$
equipped with its Borel $\sigma$-field $\mathcal E$, such that $\mathbb P$-a.s.\
the path $t\mapsto X_t$ is continuous. Let $f\colon E\to\mathbb R$ be
continuous (in particular $\mathcal E$-measurable), and set
\[
  S_t := \sup_{s\in[0,t]} f(X_s) \in (-\infty,+\infty], \qquad t\in\mathbb R_+ .
\]
\begin{enumerate}
  \item If $(\Omega,\mathcal F,\mathbb P)$ is complete, then $S_t$ is
        $\mathcal F$-measurable for every $t\in\mathbb R_+$.
  \item If $X$ is adapted to a filtration $(\mathcal F_t)_{t\in\mathbb R_+}$
        which is complete, i.e.\ $(\Omega,\mathcal F,\mathbb P)$ is complete and
        $\mathcal F_0$ contains all $\mathbb P$-negligible sets, then
        $(S_t)_{t\in\mathbb R_+}$ is adapted to $(\mathcal F_t)_{t\in\mathbb R_+}$.
\end{enumerate}
\end{proposition}

\begin{proof}
Let $\Omega'\in\mathcal F$ with $\mathbb P(\Omega')=1$ be an event on which
$t\mapsto X_t(\omega)$ is continuous. Fix $t\in\mathbb R_+$ and put
$D_t := ([0,t]\cap\mathbb Q)\cup\{t\}$, a countable dense subset of $[0,t]$.

\emph{Step 1: reduction to a countable supremum.}
Fix $\omega\in\Omega'$ and let $g(s):=f(X_s(\omega))$ for $s\in[0,t]$. As the
composition of the continuous maps $s\mapsto X_s(\omega)$ and $f$, the function
$g$ is continuous on $[0,t]$. Clearly $\sup_{D_t} g\le\sup_{[0,t]} g$.
Conversely, for $s\in[0,t]$ pick $q_n\in D_t$ with $q_n\to s$; by continuity,
\[
  g(s)=\lim_{n\to\infty} g(q_n)\le\sup_{q\in D_t} g(q).
\]
Hence, for all $\omega\in\Omega'$,
\[
  S_t(\omega)=\widetilde S_t(\omega),
  \qquad\text{where } \widetilde S_t:=\sup_{q\in D_t} f(X_q).
\]

\emph{Step 2: proof of (1).}
Each $f(X_q)$ is $\mathcal F$-measurable, since $X_q$ is
$\mathcal F/\mathcal E$-measurable and $f$ is $\mathcal E$-measurable. As a
countable supremum, $\widetilde S_t$ is $\mathcal F$-measurable with values in
$(-\infty,+\infty]$. Let $A\in\mathcal B((-\infty,+\infty])$. By Step 1,
\[
  \{S_t\in A\}
  =\bigl(\Omega'\cap\{\widetilde S_t\in A\}\bigr)
   \,\cup\,\bigl((\Omega\setminus\Omega')\cap\{S_t\in A\}\bigr).
\]
The first set lies in $\mathcal F$, because $\tilde{S}_t$ is a countable supremum of $\mc{F}$-measurable maps, making the supremum also $\mc{F}$-measurable. The second is a subset of the
$\mathbb P$-null set $\Omega\setminus\Omega'$, hence belongs to $\mathcal F$ by
completeness. Thus $\{S_t\in A\}\in\mathcal F$.

\emph{Step 3: proof of (2).}
Fix $t$. For every $q\in D_t$ we have $q\le t$, so $X_q$ is
$\mathcal F_q$-measurable and $\mathcal F_q\subseteq\mathcal F_t$; hence
$f(X_q)$, and therefore $\widetilde S_t$, is $\mathcal F_t$-measurable.
Since $\Omega\setminus\Omega'$ is $\mathbb P$-null, it lies in
$\mathcal F_0\subseteq\mathcal F_t$, so $\Omega'\in\mathcal F_t$ and
$\Omega'\cap\{\widetilde S_t\in A\}\in\mathcal F_t$. The set
$(\Omega\setminus\Omega')\cap\{S_t\in A\}$ is $\mathbb P$-negligible, hence in
$\mathcal F_0\subseteq\mathcal F_t$. By the same decomposition as in Step 2,
$\{S_t\in A\}\in\mathcal F_t$. As $t$ was arbitrary, $(S_t)$ is adapted.
\end{proof}

\begin{remark}
\leavevmode
\begin{enumerate}
  \item On $\Omega'$, $S_t$ is in fact finite: $X([0,t])(\omega)$ is compact as
        the continuous image of a compact set, so $f$ is bounded on it.
  \item If \emph{every} path is continuous, one may take $\Omega'=\Omega$ and
        completeness is not needed.
  \item Continuity of $f$ cannot simply be dropped from the proof.
        With $U\sim\mathrm{Unif}[0,1]$, $X_s=s-U$, $f=\mathbf 1_{\{0\}}$ and
        $t=1$, we have $S_1\equiv 1$, whereas
        $\sup_{q\in D_1}f(X_q)=\mathbf 1_{\{U\in\mathbb Q\}}=0$ a.s.
  \item The conclusion nevertheless remains true for Borel $f$. For
        $c\in\mathbb R$,
        \[
          \{S_t>c\}\cap\Omega'
          =\pi_\Omega\bigl(\{(s,\omega)\in[0,t]\times\Omega' :
             f(X_s(\omega))>c\}\bigr),
        \]
        the projection of a $\mathcal B([0,t])\otimes\mathcal F_t$-measurable
        set (continuous adapted processes are progressive). By the measurable
        projection theorem this projection is $\mathcal F_t$-measurable when the
        filtration is complete, although in general it is only analytic
        (Lusin--Suslin).
\end{enumerate}
\end{remark}

\subsection{Stopping times and associated \texorpdfstring{$\sigma$}{sigma}-algebras}
\begin{definition} \todo{stopping time; $\F_T$; $\F_{T+}$} \end{definition}
\begin{proposition} \todo{right-continuous filtration: $T$ s.t.\ iff $\{T<t\}\in\F_t$} \end{proposition}
\begin{proposition}[Properties] \todo{$S\le T\Rightarrow\F_S\subseteq\F_T$; $S\wedge T$, $S\vee T$; $\{S\le T\}\in\F_S\cap\F_T$; $T_n\downarrow T$} \end{proposition}
\begin{proposition} \todo{$X$ progressive $\Rightarrow$ $X_T\ind_{\{T<\infty\}}$ is $\F_T$-measurable; stopped process progressive} \end{proposition}
\begin{proposition}[Hitting times] \label{chap3.2:proposition-hitting_times} \todo{open set: s.t.\ for $\F_{t+}$; closed set + continuous paths: s.t.} \end{proposition}
\begin{quizbox} Hitting-time results (which set type needs which hypothesis) and $X_T\in\F_T$. \end{quizbox}





\subsection{Continuous-time martingales and supermartingales}
\begin{definition} \todo{(sub/super)martingale} \end{definition}
\begin{example}[Brownian martingales] \todo{$B_t$, $B_t^2-t$, $\exp(\theta B_t-\theta^2t/2)$} \end{example}
\begin{proposition}[Maximal inequality, Doob's $L^p$ inequality] \todo{via discrete version on countable dense set + right-continuity} \end{proposition}
\begin{proposition}[Upcrossing inequality] \todo{} \end{proposition}
\begin{theorem}[Regularity / càdlàg modification] \todo{right-cont.\ filtration + usual cond.\ + $t\mapsto\E X_t$ right-cont.\ $\Rightarrow$ càdlàg modification} \end{theorem}
\begin{theorem}[Convergence] \todo{a.s.\ convergence if bounded in $L^1$; closed by $X_\infty$ $\iff$ UI $\iff$ $L^1$-convergence} \end{theorem}
\begin{quizbox} Doob $L^2$: $E[\sup_{s\le t}X_s^2]\le4E[X_t^2]$. UI $\iff$ closed martingale. \end{quizbox}

\subsection{Optional stopping theorems}\label{sec:ost}

\begin{theorem}[Optional stopping, UI case]\label{chap3.4:theorem-optional_stopping_theorem_for_martingales_UI_criteria}
\todo{$S\le T$ s.t., $X$ UI càdlàg: $E[X_T\mid\F_S]=X_S$} \end{theorem}
\begin{corollary} \todo{bounded stopping times; stopped martingale $X^T$ is a martingale} \end{corollary}
\begin{proposition}[Supermartingale version] \todo{nonnegative supermartingale} \end{proposition}
\begin{example}[Applications] \todo{exit of $(-a,b)$: $\PP(T_{-a}<T_b)=\tfrac{b}{a+b}$, $\E T=ab$; $E[e^{-\lambda T_a}]=e^{-a\sqrt{2\lambda}}$} \end{example}
\begin{quizbox} Every application in the example: be fast at choosing the right martingale and justifying the limit (dominated/bounded convergence). \end{quizbox}
\begin{trap} $E[B_{T_a}]=a\ne0$: OST fails without UI/boundedness. \end{trap}
